In this subsection, we will discuss a heuristics-based implementation of a function
\[ \delta\reuse : \ty{Const} \rightarrow \ty{Fn}_{\textit{RC}} \]
that inserts \kw{reset}/\kw{reuse} instructions.
Given $\kw{let}\ z = \kw{reset}\ x$, we remark that, in every control path, $z$ may
appear at most once, and in one of the following two instructions:
$\kw{let}\ y = \kw{reuse}\ z\ \kw{ctor}_i\ \many{w}$, or $\kw{dec}\ z$.
We use $\kw{dec}\ z$ for control paths where $z$ cannot be reused.
We implement the function $\delta\reuse$ as
\[ \delta\reuse(c) = \lambda\ \many{y}.\ R(F) \text{ where } \delta(c) = \lambda\ \many{y}.\ F \]
The function $R(F)$ (\cref{fig:resetreuse}) uses a simple heuristic for replacing
$\kw{ctor}_i\ \many{y}$ expressions occurring in $F$ with
$\kw{reuse}\ w\ \kw{in ctor}_i\ \many{y}$ where $w$ is a fresh
variable introduced by $R$ as the result of a new $\kw{reset}$
operation. For each arm $F_i$ in
a $\kw{case}\ x\ \kw{of}\ \many{F}$ operation, the function $R$ requires the arity $n$
of the corresponding matched constructor. In the actual implementation,
we store this information for each arm when we compile our typed frontend language into \pureIR{}.
The auxiliary functions $D$ and $S$ implement the \emph{dead} variable search and \emph{substitution} steps respectively.
For each \kw{case} operation, $R$ attempts to insert \kw{reset}/\kw{reuse} instructions
for the variable matched by the \kw{case}. This is done using $D$ in each arm of the \kw{case}.
Function $D(z, n, F)$ takes as parameters the variable $z$ to reuse and the arity $n$ of the matched constructor.
$D$ proceeds to the first location where $z$ is dead, i.e. not used in the remaining function body,
and then uses $S$ to attempt to find and substitute a \emph{matching}
constructor $\kw{ctor}_i\ \many{y}$ instruction with a $\kw{reuse}\ w\ \kw{in}\ \kw{ctor}_i\ \many{y}$ in the remaining code.
If no matching constructor instruction can be found, $D$ does not modify the function body.
%
\begin{figure}
  \begin{flalign*}
    &R : \ty{FnBody}\pure \rightarrow \ty{FnBody}\RC \\
    &R(\kw{let}\ x = e;\ F) = \kw{let}\ x = e;\ R(F) \\
    &R(\kw{ret}\ x) = \kw{ret}\ x \\
    &R(\kw{case}\ x\ \kw{of}\ \many{F}) = \kw{case}\ x\ \kw{of}\ \many{D(x, n_i, R(F_i))} && \\
      &\phantom{==} \text{where } n_i = \#\text{fields of } x \text{ in } i\text{-th branch}
  \end{flalign*}

\begin{flalign*}
  &D : \ty{Var} \times \mathbb{N} \times \ty{FnBody}\RC \rightarrow \ty{FnBody}\RC && \\
  &D(z, n, \kw{case}\ x\ \kw{of}\ \many{F}) = \kw{case}\ x\ \kw{of}\ \many{D(z, n, F)} \\
  &D(z, n, \kw{ret}\ x) = \kw{ret}\ x \\
  &D(z, n, \kw{let}\ x = e;\ F) = \kw{let}\ x = e;\ D(z, n, F) \\
  &\phantom{==} \text{if } z \in e \text{ or } z \in F \\
  &D(z, n, F) = \kw{let}\ w = \kw{reset}\ z;\ S(w, n, F) \\
  &\phantom{==} \text{otherwise, if } S(w, n, F) \ne F \text{ for a fresh }w \\
  &D(z, n, F) = F \phantom{=} \text{otherwise}
\end{flalign*}

\begin{flalign*}
  &S : \ty{Var} \times \mathbb{N} \times \ty{FnBody}\RC \rightarrow \ty{FnBody}\RC && \\
  &S(w, n, \kw{let}\ x = \kw{ctor}_i\ \many{y};\ F) = \kw{let}\ x = \kw{reuse}\ w\ \kw{in}\ \kw{ctor}_i\ \many{y};\ F \\
  &   \phantom{==} \text{if } \abs{\many{y}} = n \\
  &S(w, n, \kw{let}\ x = e;\ F) = \kw{let}\ x = e;\ S(w, n, F) \phantom{=} \text{otherwise} \\
  &S(w, n, \kw{ret}\ x) = \kw{ret}\ x \\
  &S(w, n, \kw{case}\ x\ \kw{of}\ \many{F}) = \kw{case}\ x\ \kw{of}\ \many{S(w, n, F)}
\end{flalign*}
\caption{Inserting \kw{reset}/\kw{reuse} pairs}
\label{fig:resetreuse}
\end{figure}

As an example, consider the $\textit{map}$ function for lists
\begin{lstlisting}
  map f xs = case xs of
    (ret xs)
    (let x = proj_1 xs; let s := proj_2 xs;
     let y = f x; let ys = map f s;
     let r = ctor_2 y ys; ret r)
\end{lstlisting}
Applying $R$ to the body of $\textit{map}$ , we have $D$ looking for opportunities to \kw{reset}/\kw{reuse} $\textit{xs}$
in both \kw{case} arms.
Since $\textit{xs}$ is unused after $\kw{let}\ s = \kw{proj}_2\ \textit{xs}$, $S$ is applied to the rest of the function,
looking for constructor calls with two parameters. Indeed, such a call can be found in the let-binding for $r$. Thus,
function $D$ successfully inserts the appropriate instructions, and we obtain the function described in \cref{sec:examples}.
%% \begin{lstlisting}
%%   map f xs = case xs of
%%     (ret xs)
%%     (let x = proj_1 xs; inc x; let s := proj_2 xs; inc s;
%%      let w = reset xs;
%%      let y = f x; let ys = map f s;
%%      let r = (reuse w in ctor_2 y ys); ret r)
%% \end{lstlisting}
Now, consider the list \emph{swap} function that swaps the first two elements of a list.
It is often defined as
\begin{lstlisting}
  swap [] = []
  swap [x] = [x]
  swap (x:y:zs) = y:x:zs
\end{lstlisting}
In \pureIR{}, this function is encoded as
\begin{lstlisting}
  swap xs = case xs of
     (ret xs)
     (let t_1 = proj_2 xs; case t_1 of
       (ret xs)
       (let h_1 = proj_1 xs;
        let h_2 = proj_1 t_1; let t_2 = proj_2 t_1;
        let r_1 = ctor_2 h_1 t_2; let r_2 = ctor_2 h_2 r_1; ret r_2))
\end{lstlisting}
By applying $R$ to \emph{swap}, we obtain
\pagebreak
\begin{lstlisting}
  swap xs = case xs of
    (ret xs)
    (let t_1 = proj_2 xs; case t_1 of
      (ret xs)
      (let h_1 = proj_1 xs; let w_1 = reset xs;
       let h_2 = proj_1 t_1; let t_2 = proj_2 t_1;
       let w_2 = reset t_1; let r_1 = reuse w_2 in ctor_2 h_1 t_2;
       let r_2 = reuse w_1 in ctor_2 h_2 r_1; ret r_2))
\end{lstlisting}
Similarly to the $\textit{map}$ function, the code generated for the
function $\textit{swap}$ will \emph{not} allocate any memory when the
list value is not shared. This example demonstrates that our heuristic
procedure can avoid memory allocations even in functions containing
many nested $\kw{case}$ instructions. The example also makes it clear
that we could further optimize our \rcIR{} by adding additional
instructions. For example, we can
add an instruction that combines $\kw{reset}$ and $\kw{reuse}$ into a single instruction
and is used in situations where $\kw{reuse}$ occurs \emph{immediately}
after the corresponding $\kw{reset}$ instruction such as in the example above where we have
$\kw{let}\ w_2 = \kw{reset}\ t_1;\ \kw{let}\ r_1 =\kw{reuse}\ w_2\ \kw{in}\ \kw{ctor}_2\ h_1\ t_2$,

% For example,  where a

%% \begin{thm}
%%   For any constant $c \in \dom(\delta)$, we have
%%     \[\delta\reuse \vdash\reuse c\]
%%     and
%%     \[\erase\RC(\delta\reuse(c)) = \delta(c)\]
%% \end{thm}
%% \begin{proof}
%%   By induction over $\delta(c)$'s body.
%% \end{proof}

%% In the remaining text, instead of depending on a specific implementation of
%% $\delta\reuse$, we will merely assume these two properties.

%%% Local Variables:
%%% mode: latex
%%% TeX-master: "refcount"
%%% End:
